\section{The Euler class of a single spin operator}\label{sec:spin-residues}
A single spin line yields a real Euler class before pairing the two normal directions. This class retains local collision data which disappear from the closed partition function. We compute it first on the finite-map component and then include the stable maps with a contracted elliptic tail.
Let $C$ be a smooth closed complex curve of genus $g\ge2$, and let
$\theta$ be a theta characteristic on $C$, a holomorphic line with an
isomorphism $\theta^2\simeq K_C$. If $f:Y\to C$ is a holomorphic
cover, its pullback spin operator on $L=f^*\theta$ is the real
Cauchy--Riemann operator obtained by adding the metric conjugate of
$df:L\to K_YL^{-1}$ to $\bar\partial_L$. Lee and Parker
\cite{lpob} constructed its real obstruction bundle on the space of
stable maps. Its restriction to unramified covers has rank zero, and
the orientation of a regular cover is its spin parity. Ramification
produces a bundle of positive rank. We determine its Euler class on
the first Hurwitz spaces on which this occurs.

Write $H_m=1+m(g-1)$. In degree $m\ge2$ and source genus $H_m+1$, the
connected unmarked finite maps form a compact smooth stack surface
$\mathcal F_m$. Its general point is a cover with two simple
ramification points. For a connected boundary divisor $B$ of nodal
covers put $s_\theta(B)=(-1)^{h^0(\widetilde f^*\theta)}$, where
$\widetilde f$ is the normalization of a member. For a divisor of
covers with a single triple ramification point $p$, put
$s_\theta(B)=(-1)^{h^0(f^*\theta(p))}$. Both are parities of theta
characteristics. Our divisor formula is
\begin{equation}\label{p8-eq:intro-divisor}
       e(\operatorname{coker}D_\theta)
          =-\frac14\sum_{B} s_\theta(B)[B]
             \quad\text{in }H^2(\mathcal F_m;\QQ).
\end{equation}
Here $B$ runs over the reduced divisors of the two indicated types.
At two simple ramification points over the same value the obstruction
has rational Euler class zero. The equality holds before integrating
against another class.

\subsection{The spin obstruction and the residue pairing}\label{p8-sec:residue}

Let \(f:Y\to C\) be a finite map of smooth compact curves and fix an
isomorphism \(\theta^2\simeq K_C\). Put \(L=f^*\theta\) and
\(P=K_YL^{-1}\). The differential of \(f\) gives a holomorphic map
\(\sigma:L\to P\) whose divisor is the ramification divisor \(R\).
Write \(\deg R=2r\). A Hermitian metric on \(L\) gives a conjugate-linear
identification \(c:P\to\Lambda^{0,1}T^*Y\otimes L\). The real spin
operator and the complex space of ramification residues are
\begin{equation}\label{p8-eq:residue-operator}
 D_\theta=\bar\partial_L+c\sigma,\qquad
 \mathcal E_f=H^0(Y,P/\sigma L).
\end{equation}
The identity
\(\|D_\theta u\|^2=\|\bar\partial u\|^2+\|c\sigma u\|^2\)
follows from Stokes' theorem applied to \(\sigma u^2\). In particular
\(D_\theta\) is injective and its cokernel \(V_\theta\) has real rank
\(2r\). We orient it by continuing the complex determinant orientation
along \(\bar\partial_L+t c\sigma\), \(0\le t\le1\), with the dual
determinant convention for a cokernel \cite[\S12]{lpob}.

\begin{theorem}\label{p8-thm:residue}
Over a family of finite maps with smooth domains and fixed total
ramification \(2r\), the complexification of \(V_\theta\) identifies
with the complex orthogonal bundle \(\mathcal E\) whose form is
\begin{equation}\label{p8-eq:ramification-residues}
 b(v,v')=\sum_{p\in R_{\rm red}}\operatorname{Res}_p
                         \frac{vv'}{\sigma}.
\end{equation}
More precisely, up to the universal normalization of the Serre pairing,
\(V_\theta\) identifies with a positive real form of \(\mathcal E\).
At each map the image \(S=H^0(P)/\sigma H^0(L)\) is a maximal
isotropic subspace. Projection from \(V_\theta\) to
\(\mathcal E/S=S^*\) is a real isomorphism of orientation sign
\((-1)^{h^0(L)}\).
\end{theorem}

\begin{proof}
Resolve \([L\xrightarrow{\sigma}P]\) by its Dolbeault complex. In
degrees zero, one and two the terms are \(\Omega^0(L)\),
\(\Omega^{0,1}(L)\oplus\Omega^0(P)\) and \(\Omega^{0,1}(P)\).
The differentials, written in local coefficients, are
\[
 d_0u=(\bar\partial u,\sigma u),\qquad
 d_1(\beta,v)=\bar\partial_Pv-\sigma\beta .
\]
The sheaf map is injective and its cokernel has finite support.
Hence the middle cohomology is \(\mathcal E_f\) and all other
cohomology vanishes. On the middle term take the symmetric bilinear
pairing
\[
     B\big((\beta,v),(\beta',v')\big)
                      =\int_Y(\beta v'+\beta'v),
\]
with the one-form factors in a fixed order. Integration by parts
identifies the \(B\)-annihilator of \(\operatorname{im}d_0\) with
\(\ker d_1\). On cohomology \(B\) is \eqref{p8-eq:ramification-residues}
up to a constant. Indeed, lift local representatives of \(v\) at the
zeros of \(\sigma\) by a cutoff to a smooth section \(\widetilde v\)
of \(P\). A closed middle cochain representing \(v\) is then
\((\bar\partial\widetilde v/\sigma,\widetilde v)\). Pairing two such
lifts and applying Stokes off disks about the zeros gives the sum of
the boundary integrals of \(vv'/\sigma\). At a zero of order \(d\),
the local form on \(\CC[z]/(z^d)\) takes the coefficient of \(z^{d-1}\)
in the product, after multiplication by a unit. It is nondegenerate,
and distinct support points are orthogonal.

The conjugate-linear involution
\[
                J(\beta,v)=(c(v),c^{-1}(\beta))
\]
makes \(B(\,\cdot\,,J\,\cdot\,)\) a positive Hermitian \(L^2\)
metric after a constant normalization. In a holomorphic frame of
squared norm \(h\), \(J\) sends the local coefficients \((\beta,v)\)
to \((h^{-1}\overline v,h\overline\beta)\). The harmonic middle
space is the Hermitian orthogonal complement of
\(\operatorname{im}d_0+J\operatorname{im}d_0\). In particular it is
\(J\)-invariant. Projection to the first factor identifies the
\(J\)-fixed part of the middle term with \(\Omega^{0,1}(L)\), and
sends \((1+J)d_0u/2\) to \(D_\theta u/2\). A fixed vector real
orthogonal to all these projections is Hermitian orthogonal to
\(\operatorname{im}d_0\), by testing both \(u\) and \(iu\). It is
then also orthogonal to \(J\operatorname{im}d_0\). Thus the real
cokernel identifies with the fixed harmonic space. Elliptic
stabilization gives the identification in families since the harmonic
dimension is constant. Equivalently the orthogonal structure group of
\(\mathcal E\) reduces from \(\mathrm O(2r,\CC)\) to
\(\mathrm O(2r,\RR)\).

The residue theorem makes \(S\) isotropic, and Riemann--Roch gives
\(\dim_\CC S=r\). The long exact sequence identifies
\(\mathcal E_f/S\) with the annihilator of \(\sigma H^0(L)\) in
\(H^1(L)=H^0(P)^*\). On a harmonic class it is the map
\((\beta,v)\mapsto[\beta]\). The image of the zeroth-order term \(c\sigma\) on the Dolbeault kernel is the Hermitian dual of \(\sigma H^0(L)\) in \(H^1(L)\).
Consequently the map from the fixed harmonic space to \(S^*\)
is the cokernel projection of Lee and Parker
\cite[Lemma~12.1 and Proposition~12.2]{lpob}. Its sign is
\((-1)^{h^0(L)}\), from conjugating the complex kernel before taking
its determinant. This last comparison is at a single map and uses
no constant-rank assumption on the Dolbeault cohomology of nearby maps.
\end{proof}

Write \(\kappa=c_1(K_C)\) and use rational cohomology. On a stratum
of fixed ramification profile we may pass to a finite cover ordering
the ramification points. Denote by \(\operatorname{ev}_p\) evaluation
at a marked ramification point \(p\). The following formula concerns
the restriction to this stratum.

\begin{theorem}\label{p8-thm:profiles}
If a ramification index is even, \(e(V_\theta)=0\). If all indices
are odd, write \(R=2D\), \(A=L(D)\), and put
\(I=H^0(Y,L(D)/L)\). Then
\begin{equation}\label{p8-eq:odd-profile-euler}
                 e(V_\theta)=(-1)^{h^0(A)}c_r(I).
\end{equation}
For a fixed target an index \(2s+1\ge5\) also gives zero. If the
ramification consists of \(r\) points of index three, the restriction
to the ordered stratum is
\begin{equation}\label{p8-eq:triple-profile-euler}
             e(V_\theta)=(-1)^{h^0(A)}6^{-r}
                           \prod_{p\in R_{\rm red}}\operatorname{ev}_p^*\kappa .
\end{equation}
\end{theorem}

\begin{proof}
At an even index the corresponding orthogonal residue summand has odd
complex rank. A sum of complex orthogonal bundles admits a positive
real form obtained by taking a positive real form in each summand.
Any two such reductions are isomorphic, since the space of reductions
in each fiber is contractible. After orienting the summands on a
finite cover, the Euler class thus has the Euler class of an
odd-rank oriented real bundle as a factor. The latter vanishes over
\(\QQ\). Transfer gives the assertion on the original stratum.

For odd indices \(A^2=K_Y\) by Riemann--Hurwitz. Multiplication by
\(\sigma\) identifies \(P\) with \(L(2D)\). In this description \(I\)
is a maximal isotropic subbundle of \(\mathcal E\), since the local
meromorphic representatives in \(I\) have product with no residue
in \eqref{p8-eq:ramification-residues}. The two maximal isotropics have
\[
                     k=\dim(S\cap I)=h^0(A)-h^0(L).
\]
This is the residue construction for a theta characteristic
\cite[\S1]{mumford}. The line \(A\) is the theta characteristic of the
oddly ramified cover used in spin Hurwitz theory \cite[\S1]{lp}.
For two maximal isotropics meeting in dimension \(k\), the real
projections from a positive real form to their duals induce
orientations which differ by \((-1)^{r-k}\). To see this, move the
pair through maximal isotropics with intersection of dimension \(k\)
to a pair in hyperbolic coordinate form. In the \(r-k\) exchanged
coordinates the projections differ by complex conjugation. Applying
Theorem~\ref{p8-thm:residue} and \(c_r(I^*)=(-1)^r c_r(I)\) gives the
sign \((-1)^{h^0(L)+r-k+r}=(-1)^{h^0(A)}\). The bundle \(I\) exists
over the entire profile stratum. The parity of the theta
characteristic \(A\) is locally constant, so the pointwise orientation
comparison proves \eqref{p8-eq:odd-profile-euler} across jumps of \(h^0(L)\).

At an index \(2s+1\) point with image \(x\), the filtration by pole
order on \(I\) has graded lines
\(\theta_x\otimes(T_pY)^{\otimes j}\), \(1\le j\le s\). The
leading nonzero jet of \(f\) identifies \((T_pY)^{\otimes(2s+1)}\)
with \(T_xC\). The first Chern classes of the graded lines are
therefore
\[
              \left(\frac12-\frac{j}{2s+1}\right)
                         \operatorname{ev}_p^*\kappa .
\]
When \(s\ge2\), two such factors already have product zero since
\(\kappa^2=0\) on \(C\). For \(s=1\) the factor is
\(\operatorname{ev}_p^*\kappa/6\), giving
\eqref{p8-eq:triple-profile-euler}.
\end{proof}

On the stratum with \(r\ge2\) distinct triple critical values,
the product in \eqref{p8-eq:triple-profile-euler} itself vanishes in
ordinary cohomology. It pulls back from the top cohomology of the
noncompact ordered configuration space of \(r\) points on \(C\).
A prescribed coincidence of two values gives zero by \(\kappa^2=0\).
The class of a family meeting the boundary can nevertheless have a
contribution supported there.

\begin{corollary}\label{p8-cor:nodal-spin-matching}
On a stratum of finite nodal maps with fixed dual graph, label the
components, the nodes, and the two preimages of each node on a finite
cover. Let \(V_{\theta,i}\) be the oriented cokernel on the \(i\)-th
normalized component and let \(\operatorname{ev}_j\) send a map to
the image \(x_j\) of its \(j\)-th node.
Then the restriction of the nodal Euler class is
\[
             \prod_i e(V_{\theta,i})\,
             \prod_j\operatorname{ev}_j^*c_1(\theta).
\]
An unramified component contributes its index-zero determinant sign.
\end{corollary}

\begin{proof}
In the \(W^{1,p}\) domain with \(p>2\), the nodal domain is the
kernel of the surjection taking the differences of values over the
nodes, with target \(\bigoplus_j\theta_{x_j}\). The spin operators
on the normalized components are injective. The exact sequence of
kernels and cokernels therefore gives
\[
 0\longrightarrow\bigoplus_j\theta_{x_j}\longrightarrow
 V_{\theta,\mathrm{nod}}\longrightarrow\bigoplus_iV_{\theta,i}
 \longrightarrow0.
\]
The determinant isomorphism gives the matching lines their complex
orientations, by continuation from the complex-linear normalization
sequence \cite{zinger}. All ranks in the sequence are even, so
reordering the factors preserves the product orientation. Interchange
of the two branches at a node tensors its matching line with a flat
sign line and leaves its rational first Chern class unchanged.
\end{proof}
